% ------------------------------------------------------------
% §5  Frameworks for analysis, and the gap
% ------------------------------------------------------------
\section{Frameworks, and the size of the gap}\label{sec:frameworks}

We now assemble the existing theoretical apparatus around \SHA{} and take
honest stock of how far it reaches. The summary in advance: there are
mature frameworks on \emph{both} sides of the object---idealized models
above it, empirical and cryptanalytic measurements below it---and a
canyon in the middle where the function itself lives.

\subsection{What ``pseudorandom'' officially means}
\label{sec:definitions}

The foundational definitions of pseudorandomness are about
\emph{families}, not single objects. A pseudorandom function family in
the sense of Goldreich, Goldwasser, and Micali~\citep{GGM1986} is a keyed
collection $\{F_k\}$ such that no efficient algorithm, allowed to query a
black box, can tell a randomly keyed $F_k$ from a truly random function
with nonnegligible advantage. This definition is a triumph---it supports
composition theorems, equivalences, a whole complexity theory of
randomness---but notice that it \emph{cannot even be stated} for the
fixed, keyless, public function \SHA{}. For any fixed function, a trivial
\emph{distinguisher} exists. A distinguisher is the formal adversary in
these definitions: an efficient algorithm that queries a black box
containing either the function under test or a truly random function,
emits a one-bit verdict, and is scored by its \emph{advantage}, the gap
between its acceptance probabilities in the two worlds. The notion
descends from Yao's theory of computational
indistinguishability~\citep{Yao1982} and receives its textbook treatment
in Chapter~3 of~\citet{Goldreich2001}. In our context the trivial
distinguisher simply hard-codes one known value: query the box at
$\mathsf{0}$ and accept if the answer is the published
$\SHA(\mathsf{0})$. Against \SHA{} it accepts always; against a random
function, with probability $2^{-256}$; its advantage is essentially $1$,
and no keyless public function escapes it. The same
foundational awkwardness afflicts collision resistance: since collisions
exist by pigeonhole, ``no efficient algorithm finds one'' is technically
false---some two-line program prints a collision; humanity just cannot
exhibit it. Rogaway's ``human ignorance'' formalism~\citep{Rogaway2006}
repairs the semantics by making theorems constructive (explicit
reductions transforming any breaker of a protocol into an explicit
collision-finder), and it is the honest logical home of every security
claim about \SHA{}. But note what has happened: \emph{the definition of
success has been moved from mathematics into history}---``no one has
written down a collision, so far.''

\subsection{Idealized models: reasoning above the object}
\label{sec:models}

Protocol designers escape upward. In the random-oracle
model~\citep{BellareRogaway1993} one simply \emph{postulates} a public
uniformly random function and proves protocols secure relative to it; the
leap back to \SHA{} is heuristic, frankly acknowledged, and empirically
successful. The model even has an internal quality-control theory:
\emph{indifferentiability}~\citep{MRH2004} asks whether a construction
built from an ideal component is safely interchangeable with an ideal
whole, and by this standard plain Merkle--Damg{\aa}rd \emph{fails}---the
length-extension property (from $\SHA(m)$ anyone computes
$\SHA(m \Vert \mathrm{pad} \Vert m')$ without knowing $m$) is a genuine,
exploitable deviation from random-oracle behavior, identified precisely
and repaired by variants with security proofs~\citep{CDMP2005}. This is
the theory at its best: it cannot certify $f$ itself, but given an
idealized $f$ it classifies \emph{constructions} completely. The
architecture above the compression function is, in a real sense, solved
mathematics---the two-layer picture of \cref{sec:constructions}, where
the Merkle--Damg{\aa}rd mode (\cref{thm:md}) and the Davies--Meyer
compression (\cref{thm:dm}) each carry a theorem---and all residual
mystery is compressed into the finite map $f$.

There is a sharper way to say that the model is a heuristic, and it is a
theorem rather than a caveat. Canetti, Goldreich, and Halevi
constructed signature and encryption schemes that are \emph{provably
secure} in the random-oracle model and \emph{provably insecure} the
moment the oracle is replaced by \emph{any} concrete, publicly specified
function whatsoever, \SHA{} or otherwise~\citep{CGH2004}. The
construction is diagonal---the scheme tests whether the code it is given
is a short program and betrays its secret if so, which a true random
oracle never is but every real hash is---so it is contrived, and no
natural protocol is known to suffer. But it settles the logic: a
random-oracle proof is not a proof about \SHA{} with an admitted gap; it
is a proof about an idealized object that \emph{no} function can fully
become. The heuristic is not merely unproven, it is in general false,
and its continued empirical success (\cref{sec:below}) is therefore
itself part of the natural history---a regularity of the world, not a
lemma.

\subsection{The construction is not the oracle: provable structure
above \texorpdfstring{$f$}{f}}\label{sec:mdstructure}

Here is a genuine surprise, and it cuts against the drumbeat of this
paper. We keep asserting that \SHA{} exhibits \emph{no} observable
deviation from a random function. That is true of the compression
function $f$; it is \emph{false} of the Merkle--Damg{\aa}rd wrapper
around it. The iteration mode has provable, nameable, exploitable
structure---deviations from random-oracle behavior that hold
\emph{even if $f$ is a perfect random oracle}---and they are among the
most elegant results in the subject. They locate the paper's mystery
more precisely than we have so far: the residue of the unknown lives in
$f$, because the part outside $f$ is not unknown at all. Its
non-randomness is a set of theorems.

The keystone is Joux's multicollision attack~\citep{Joux2004}
(\cref{fig:multicoll}). For a true $256$-bit random function, finding
$2^k$ messages that all share one digest should cost about
$2^{256(2^k-1)/2^k}$ work---for eight-way agreement, near $2^{224}$,
astronomically beyond a single collision's $2^{128}$. In an iterated
hash it costs only $k$ times a single collision. The trick is pure
construction-level reasoning: birthday-find a colliding \emph{block
pair} at the first chaining value, then \emph{from the merged state}
find another colliding block pair, and so on for $k$ stages; freely
mixing the two choices at each stage yields $2^k$ full messages all
landing on the same final digest, at total cost $k \cdot 2^{128}$
instead of $2^{128 \cdot (2^k-1)/2^k}$. Cascading is the immediate
casualty: hashing with $H_1 \Vert H_2$, the ``belt and braces'' habit
of concatenating two hashes for safety, is no stronger than the better
single one---build a $2^{n_1/2}$-multicollision in $H_1$, then by
pigeonhole two of its exponentially many members collide in $H_2$ as
well. A folklore engineering practice was refuted by one structural
idea, no cryptanalysis of any $f$ required.

\begin{figure}[ht]
  \centering
  \begin{tikzpicture}[font=\small, >={Stealth[length=1.7mm]},
      cv/.style={circle, draw, minimum size=6.5mm, inner sep=0pt,
                 fill=projfill, font=\scriptsize},
      lbl/.style={font=\scriptsize}]
    \node[cv] (h0) at (0,0) {$h_0$};
    \node[cv] (h1) at (2.6,0) {$h_1$};
    \node[cv] (h2) at (5.2,0) {$h_2$};
    \node[cv] (h3) at (7.8,0) {$h_3$};
    \node[lbl] at (9.1,0) {$\cdots$};
    % each stage: two parallel colliding blocks
    \foreach \a/\b/\i in {h0/h1/1, h1/h2/2, h2/h3/3} {
      \draw[->] (\a) to[bend left=42]
        node[above, lbl] {$b_{\i}$} (\b);
      \draw[->] (\a) to[bend right=42]
        node[below, lbl] {$b_{\i}'$} (\b);
    }
    \node[lbl, align=center, projframe] at (3.9,-1.55)
      {each stage: one birthday collision $f(h,b_i)=f(h,b_i')$, cost
       $2^{128}$};
    \node[lbl, align=center] at (3.9,1.7)
      {choose $b_i$ \emph{or} $b_i'$ freely at every stage
       $\Rightarrow$ $2^{k}$ messages, one digest};
    \draw[decorate, decoration={brace, amplitude=4pt}]
      (0,2.05) -- (7.8,2.05);
    \node[lbl] at (10.4,0) {total $k\cdot 2^{128}$};
  \end{tikzpicture}
  \caption{Joux multicollisions~\citep{Joux2004}. Because the state is
  a narrow \emph{delay line} passed from block to block
  (\cref{sec:roundshape}), a chain of $k$ independent single collisions
  multiplies into $2^k$ colliding messages at only $k$ times the cost
  of one---a provable, non-statistical deviation from random-oracle
  behavior that lives entirely in the iteration, not in $f$. The
  narrow internal state is the culprit; a wide-pipe or sponge design
  (the SHA-3 sponge of \cref{sec:merkledamgard}) closes exactly this
  door.}
  \label{fig:multicoll}
\end{figure}

Two further theorems mine the same seam, and both convert the
\emph{narrowness of the internal state}---the eight-word delay line of
\cref{sec:roundshape}, no wider than the output---into an attack.
Kelsey and Schneier showed that a \emph{second preimage} for a long
message is far cheaper than the ideal $2^{256}$: given a target message
of $2^\ell$ blocks, ``expandable messages'' (a multicollision between
message lengths) let one splice in a forged prefix that rejoins the
victim's chain at some interior state, at cost about
$2^{256-\ell}$~\citep{KelseySchneier2005}. Long messages are
\emph{provably} weaker than short ones---a sentence with no meaning at
all for a random oracle, where every input is independent.
Kelsey and Kohno pushed the genre from collision to \emph{commitment
fraud}: their ``herding'' attack precomputes a branching ``diamond''
structure of chaining values, publishes a single digest as a
soothsayer's sealed prophecy, and then, given any later-revealed
prefix, herds it into the committed digest by appending a short
suffix~\citep{KelseyKohno2006}. The hash's advertised binding
property---the entire basis of the commitment use of
\cref{sec:uses}---is quantitatively weakened by construction-level
structure alone.

The moral reframes the whole paper, so we state it flatly. Every one of
these attacks treats $f$ as an ideal random oracle and still breaks the
advertised behavior of \SHA{}-the-construction; each is a
\emph{theorem}, not a measurement, and each traces to one visible
design fact---the internal state is no wider than the digest, so the
chaining value is a bottleneck that a random function would not have
(\cref{sec:roundshape}). This is why the SHA-3 competition selected a
sponge with a \emph{wide} internal
state~\citep{BDPV2007} that structurally forecloses all three attacks,
and why the
honest slogan is not ``\SHA{} behaves like a random oracle'' but the
sharper \emph{``$f$ behaves like a random function, and the wrapper's
deviations from a random oracle are completely known.''} The mystery of
this paper is now pinned exactly where it belongs: not in the
architecture, which is theorems all the way down, but in the single
finite map $f$, about which there are none.

\subsection{Empirical and adversarial measurement: attacking from below}
\label{sec:below}

From below, two very different measurement cultures apply.

\emph{Statistical testing} treats the function as a black-box PRNG and
runs batteries of hypothesis tests: NIST's SP\,800-22
suite~\citep{NIST2010} and the far more demanding TestU01
``BigCrush''~\citep{LEcuyerSimard2007}. \SHA{}-based generators pass
everything, which is informative but bounded: these batteries famously
demolish the classical linear generators (a Mersenne Twister fails
BigCrush's linear-complexity tests; lattice structure betrays linear
congruential generators---the pathologies catalogued since
Knuth~\citep{Knuth1997}), so passing places \SHA{} strictly above the
traditional PRNG canon; but a battery is a fixed, finite set of null
hypotheses, and structure invisible to all of them may be gross. Knuth's
own moral stands: statistical testing can refute, never confirm. For
\SHA{} the concrete record is easily stated: the counter sequence
$\SHA(1), \SHA(2), \SHA(3), \dots$ is not merely testable but
\emph{standardized}---NIST's Hash\_DRBG is essentially this
generator~\citep{SP80090A}---and the same construction over SHA-1
passed every one of BigCrush's tests in the paper that introduced the
battery. Passes of this kind are so thoroughly expected that the
literature scarcely bothers to record them. To see what that
expectation is calibrated against---and what a \emph{finished} theory
of pseudorandomness looks like---we pause (\cref{sec:prngs}).

The record deserves tabulating battery by battery, because the honest
picture distinguishes three things the phrase ``passes every test''
usually blurs: what has actually been \emph{run} and published, what is
merely \emph{expected} and therefore never written down, and what is
\emph{provably out of reach} of any feasible test at all
(\cref{tab:batteries}).

\begin{table}[H]
  \centering\small
  \renewcommand{\arraystretch}{1.25}
  \begin{tabular}{@{}p{0.15\linewidth} p{0.235\linewidth}
                     p{0.35\linewidth} p{0.115\linewidth}@{}}
    \toprule
    \textbf{Battery (size)} & \textbf{What it would catch} &
    \textbf{Status on the SHA family} & \textbf{Source} \\
    \midrule
    NIST SP\,800-22 \newline (15 tests) &
    monobit \& block frequency, runs, matrix rank, spectral, template
    matching, Maurer universal, linear complexity, serial, approximate
    entropy, cusums, excursions &
    The standard instrument for cryptographic RNGs. The \SHA{} counter
    generator is standardized as \textsf{Hash\_DRBG} and passes; a
    failure would be reportable and none is reported. Passing is
    \emph{expected}, so rarely published as such. &
    \citet{NIST2010}; \citet{SP80090A} \\
    TestU01 Crush / BigCrush \newline (${\sim}96$/$106$ tests, $160$ stats) &
    birthday spacings, collisions, gaps, poker, rank, linear complexity,
    Hamming-weight, random-walk statistics &
    \textbf{Documented run.} SHA-1 (OFB and CTR modes) passes
    \emph{all}; AES likewise, bar one $p$-value that cleared on
    replication. \SHA{} itself was not in that study, but is expected to
    behave identically. &
    \citet{LEcuyerSimard2007} \\
    DIEHARD / Dieharder \newline (${\sim}30$ tests) &
    birthday spacings, overlapping permutations, binary rank, ``monkey''
    (bitstream) tests &
    Routinely applied to cryptographic generators and passed; we know of
    no \SHA{}-specific published study---passing is folklore, exactly as
    the text describes. &
    Marsaglia \newline (software) \\
    PractRand \newline (streams to ${\gtrsim}$\,TB) &
    bit-pattern and transition anomalies visible only at very large
    sample volume &
    Applied to \SHA{}-based counter generators and passed to the volumes
    tested. A software tool, with no formal paper to cite. &
    --- (software) \\
    \addlinespace[1pt]
    \emph{Global $256$-bit uniformity / closeness} &
    \emph{any} constant-$\varepsilon$ deviation of the \emph{full} output
    distribution from uniform &
    \textbf{Not feasible.} Requires ${\sim}2^{128}$ samples---the
    birthday wall as an information-theoretic lower bound; provably
    undecidable at any affordable cost (\cref{sec:disttest}). &
    \citet{Canonne2020} \\
    \bottomrule
  \end{tabular}
  \caption{Standard statistical batteries and their verdict on the
  \SHA{} family. Only the TestU01 row is a \emph{documented, published}
  run---and it tested SHA-1, not \SHA{}. The SP\,800-22 row rests on
  standardization and the absence of reported failure; the DIEHARD and
  PractRand rows on routine practice that the literature, expecting
  success, does not record; and the last row is not a gap in effort but
  a theorem (\cref{sec:disttest}) that no feasible sample can close.
  Every battery that \emph{has} been run reports uniform; the one that
  would see the whole distribution cannot be run.}
  \label{tab:batteries}
\end{table}

It is worth being concrete about \emph{which} regularities are absent
and which are present, since ``passes the batteries'' is too coarse to
convey the texture of the thing. A ledger.

\emph{Regularities \SHA{} does \textbf{not} exhibit} (all empirical, none
proven absent):
\begin{itemize}
  \item \textbf{No lattice structure.} Successive output tuples show no
  Marsaglia effect: where the $d$-tuples of a linear congruential
  generator lie on at most $(d!\,m)^{1/d}$ hyperplanes
  (\cref{sec:prngs};~\citealp{Marsaglia1968}), the points
  $(\SHA(n), \SHA(n{+}1), \dots)$ display no detectable hyperplane
  concentration, no lattice, no low-dimensional alignment, in any
  projection anyone has computed.
  \item \textbf{No observable period.} The counter stream $\SHA(1),
  \SHA(2), \dots$ has no detected cycle; viewed as iteration
  (\cref{sec:randommaps}) the expected recurrence time is $\sim 2^{128}$,
  operationally unreachable. Contrast the classical generators, whose
  periods are exact theorems ($2^{19937}-1$ for the Mersenne
  Twister;~\citealp{MatsumotoNishimura1998}).
  \item \textbf{No linear or low-degree structure.} The $\Ftwo$-linear
  complexity of the output stream is empirically near-maximal---no low
  linear complexity, the exact defect that sinks the Twister on
  BigCrush---and no correlation of any output bit with a
  low-degree Boolean function of the input has been found. Each output
  bit, as a function of the input, appears to have a flat
  Walsh--Hadamard spectrum: no good linear approximation, the property
  linear cryptanalysis would exploit.
  \item \textbf{No first-order statistical bias.} Output bit frequencies,
  runs, gaps, block frequencies, and autocorrelations at every measured
  lag sit within sampling noise of the uniform expectation---the content
  of ``passes SP\,800-22 and BigCrush.''
\end{itemize}

\emph{Regularities \SHA{} \textbf{does} exhibit} (these are the true
structure, and they are exactly the harmless kind):
\begin{itemize}
  \item \textbf{Determinism.} It is a fixed function: $\SHA(x)$ is always
  the same, the one perfectly reliable ``pattern,'' and the sole basis of
  every use in \cref{sec:uses}.
  \item \textbf{The avalanche property.} Flipping a single input bit
  flips each of the $256$ output bits with probability indistinguishable
  from $\tfrac12$, and near-independently, so the number of changed
  output bits is essentially $\mathrm{Binomial}(256, \tfrac12)$, sharply
  peaked at $128$. This \emph{strict avalanche criterion}, a deliberate
  design target since~\citet{WebsterTavares1985}, is a regularity in the
  sense that it holds tightly and measurably---a property of
  \emph{structurelessness itself}, and the subject of a project below.
  \item \textbf{Completeness/diffusion.} After the full round count every
  output bit depends on every input bit; there is no input coordinate any
  output ignores.
\end{itemize}

The avalanche property is the one to dwell on, because it dictates the
self-correlation structure---or rather its total absence. Consecutive
counter values $n$ and $n+1$ differ, as bit strings, in just a few
low-order bits; avalanche then forces $\SHA(n)$ and $\SHA(n+1)$ to differ
in about half of \emph{their} bits, with no residual correlation between
them. So the output stream has essentially zero autocorrelation at every
lag and no continuity in any metric on the inputs: nearness of arguments
carries no information about nearness of values. This is precisely the
opposite of the individual ARX primitives, each of which is $2$-adically
$1$-Lipschitz---continuous, correlation-preserving
(\cref{sec:statespace})---and it is the composition of $64$ rounds that
annihilates that continuity. The designed regularity (avalanche) is the
mechanism that manufactures the absence of every other regularity a
statistical test could name.

\emph{Cryptanalysis} is the adversarial culture: bespoke,
design-specific, and quantitative in a different way---its currency is
the number of rounds that can be broken and at what cost.
\Cref{sec:differential} anatomizes the method behind these numbers; the
headline state of the art on \SHA{}: practical collisions for the
compression function reduced to 28 of its 64 rounds and theoretical
ones at 31~\citep{MendelNadSchlaffer2013}, ``semi-free-start''
collisions (the attacker also chooses the chaining value) at
39~\citep{LiLiuWang2024}, and preimage attacks via bicliques reaching
about 45 rounds at cost barely below brute force~\citep{KRS2012}. Two readings of these numbers
coexist. Comforting: after twenty years of world-class effort, half the
rounds stand untouched, and nothing approaches the full function.
Discomfiting: the same community's tools went, for SHA-1, from academic
paper to industrial collision in one decade (\cref{sec:falling}), and
there is no theory predicting which trajectory \SHA{} is on. The round
margin is a measurement of \emph{our} progress, not of \emph{its}
strength.

\subsection{How much can statistics even see? The sublinear ceiling}
\label{sec:disttest}

Knuth's moral---testing refutes, never confirms---has a quantitative
refinement worth stating, because it says precisely \emph{which}
questions about \SHA{}'s output distribution are answerable with a
feasible number of samples and which are not. The relevant theory is
\emph{distribution testing}, the sublinear-algorithms study of how many
independent samples from an unknown distribution over a domain of size
$N$ are needed to decide a property of it; the modern survey is
\citet{Canonne2020}.

The foundational result is the sample complexity of \emph{uniformity
testing}: to distinguish ``the distribution is uniform on $N$ points''
from ``it is $\varepsilon$-far from uniform in total-variation
distance,'' one needs, and suffices with, $\Theta(\sqrt{N}/\varepsilon^{2})$
samples. The upper bound is achieved by a startlingly simple statistic---%
\emph{count collisions}: draw $m$ samples, count coinciding pairs, and
compare against the $\binom{m}{2}/N$ expected under uniformity, a test
going back to Goldreich and Ron and made sample-optimal by Paninski
(both in~\citealp{Canonne2020}). The reader will recognize the birthday
calculus of \cref{sec:randommaps}: collisions among $\sqrt{N}$ samples
are exactly the signal, and $\sqrt{N}$ is exactly the threshold.

Point this at \SHA{}'s $256$-bit output distribution ($N = 2^{256}$) and
the ceiling is stark: detecting \emph{any} constant-$\varepsilon$
deviation of the global output distribution from uniform requires on the
order of $2^{128}$ samples---the birthday wall again, now wearing the
hat of an \emph{information-theoretic lower bound} on testing rather than
an algorithmic cost. No battery, however clever, that draws fewer than
$\sim 2^{128}$ outputs can certify or refute uniformity of the full
distribution; the question is not hard, it is \emph{unanswerable} at
feasible sample sizes. So whatever SP\,800-22 and BigCrush are doing when
they ``pass \SHA{},'' they are emphatically \emph{not} testing the
$256$-bit distribution.

What they \emph{are} testing is the reconciling point, and it is exactly
the sublinear-feasible part. Every test in those batteries is a
low-dimensional or low-complexity statistic---the frequency of single
bits, of short blocks, of runs; linear complexity; a fixed template
count---each of which lives on a domain small enough ($2^{k}$ for modest
$k$, or a real-valued summary) that $\sqrt{\,\cdot\,}$ samples suffice
and a laptop has them. Passing means: on every \emph{marginal} anyone has
thought to compute and can afford to test, \SHA{} matches uniform. The
$2^{128}$ ceiling says the joint, high-dimensional distribution is
untestable; the batteries probe its cheap shadows, and the shadows are
clean. That precise divorce---feasible low-dimensional marginals all
uniform, infeasible global distribution forever unknown---is the honest
content of ``passes every test,'' and it is a sharper statement than the
phrase usually carries.

Two further refinements are worth flagging because each squeezes signal
from limited samples rather than brute count. First, \emph{closeness}
(two-sample) testing: deciding whether \SHA{}'s output distribution
\emph{equals} a reference distribution, rather than testing it against
the explicit uniform, has its own sublinear theory~\citep{Batu2013}, and
it is the right frame for the comparative program of \cref{sec:randommaps}---%
one does not test \SHA{} against ``uniform'' in the abstract but against
a measured random-map \emph{control}, in the low-dimensional Flajolet
coordinates where $\sqrt{\,\cdot\,}$ samples reach. Second, the
\emph{$p$-value-of-$p$-values} trick that SP\,800-22 actually
uses~\citep{NIST2010}: run one weak test on many independent streams,
then test whether the resulting $p$-values are themselves uniform on
$[0,1]$. A bias too faint for any single run to flag becomes visible in
the second-order distribution of the $p$-values---a genuine way to
amplify a subtle distributional signal without paying for
prohibitively many samples per test. Against \SHA{} these meta-tests, too,
report uniform. None of this proves anything; but it locates, to the
sample, the boundary between what the empirical program can decide and
what it provably cannot.

\begin{project}{Collision testing at the birthday threshold}{probability;
  programming}{4--6 weeks}
  Implement the collision-count uniformity tester on the $n$-bit output
  of truncated \SHA{} for $n$ up to ${\sim}48$, and measure its
  \emph{sample complexity} empirically: how many outputs are needed to
  reject uniformity for reduced-round or deliberately-biased variants,
  and how does the number scale with $n$? Confirm the
  $\Theta(\sqrt{2^{n}})$ law~\citep{Canonne2020} and watch it become the
  same $2^{n/2}$ wall the collision-hunting project of
  \cref{sec:telescope} runs into---one barrier, met from the
  testing side and the attacking side at once. Then add the
  $p$-value-of-$p$-values layer and see how much fainter a bias it can
  catch at fixed budget.
\end{project}

\subsection{Interlude: what a finished theory looks like}\label{sec:prngs}

Two mature bodies of mathematics sit directly adjacent to our object,
and both show what ``understanding a pseudorandom process'' means when
understanding is actually achieved.

\paragraph{Pseudorandom numbers and finite fields.}
Pseudorandom number generators were born with the stored-program
computer: Monte Carlo simulation demanded rivers of reproducible random
digits, physical sources were slow and unrepeatable, and by mid-century
von Neumann's middle-square method and Lehmer's linear congruential
generator $x_{n+1} = a x_n + c \bmod m$ had founded the subject (von
Neumann also supplied its founding joke, about the ``state of sin'' of
anyone producing random digits by arithmetic). The subject's demands
are precise and \emph{non-adversarial}: a simulation needs
equidistribution of the statistics it consumes, a long period, and
above all speed---a physics run may consume $10^{12}$ variates, so a
generator must cost a few word operations per output. The discipline
was codified by \citet[Ch.~3]{Knuth1997}, and its mature mathematical
home is the arithmetic of finite fields, as systematized by Lidl and
Niederreiter~\citep{LidlNiederreiter1997, Niederreiter1992}. And here
the theorems are \emph{theorems}. A congruential generator attains full
period exactly when three elementary congruence conditions on
$(a, c, m)$ hold \citep[the Hull--Dobell criterion;][\S 3.2.1]{Knuth1997}.
A linear recurrence over $\Ftwo$ with characteristic polynomial $f$ has
period equal to the order of $f$, maximal $2^k - 1$ precisely when $f$
is primitive~\citep[Ch.~8]{LidlNiederreiter1997}---which is how
Matsumoto and Nishimura could \emph{prove} that their Mersenne Twister
has period $2^{19937}-1$ and equidistribution in $623$
dimensions~\citep{MatsumotoNishimura1998}. Marsaglia \emph{proved} that
successive $d$-tuples from any congruential generator fall on at most
$(d!\,m)^{1/d}$ hyperplanes~\citep{Marsaglia1968}---for $d = 10$ and
$m = 2^{32}$, a few dozen planes carrying every point the generator
will ever produce. For nonlinear congruential generators,
equidistribution is controlled through exponential-sum estimates, so
that Weil's Riemann hypothesis for curves over finite fields quietly
underwrites the uniformity of practical
generators~\citep{Niederreiter1992}. Classical pseudorandomness, in
short, is a satellite of the arithmetic of finite fields.

Now the double-edged contrast. These theorems exist \emph{because} the
generators are algebraically transparent---one clean structure, fully
exposed---and the same transparency is what makes them cryptographically
worthless: Marsaglia's planes and the Twister's $\Ftwo$-linearity are
simultaneously the content of the theorems and the substance of the
distinguishers (\cref{sec:below}). Efficiency tells the story from the
cost side. A Twister step is a handful of word operations; \SHA{}
spends on the order of a few thousand word operations per $256$-bit
output---one to two orders of magnitude more work per bit in
software---and the entire surcharge purchases exactly one property,
resistance to an intelligent adversary, which no simulation needs.
Where the classical theory lives inside a single algebraic structure,
\SHA{} deliberately straddles two incompatible ones
(\cref{sec:statespace}) precisely so that no analogue of the
Lidl--Niederreiter machinery can grip it. That is the trade: the
theorems stop where the adversary begins.

\paragraph{Random permutations and the symmetric group.}
The second finished theory concerns a different pseudorandom object: a
\emph{shuffle} is a pseudorandom number generator valued in the
symmetric group, and ``how many riffle shuffles randomize a deck?'' is
the question ``how many rounds suffice?'' asked of cards. Here the
theory reached a summit. Modeling repeated riffle shuffles as a random
walk on $S_{52}$, \citet{BayerDiaconis1992} computed the total-variation
distance to uniformity essentially exactly, exhibiting the famous
answer (seven shuffles) and, more importantly, the \emph{cutoff
phenomenon}. Their theorem is worth quoting in the exact form that makes
the abruptness visible. Riffle-shuffling a deck of $n$ cards
$k = \tfrac{3}{2}\log_2 n + \theta$ times, the total-variation distance
to the uniform distribution obeys
\[
  \bigl\| Q^{\ast k} - U \bigr\|_{\mathrm{TV}}
  \;=\; 1 - 2\,\Phi\!\Bigl(\tfrac{-2^{-\theta}}{4\sqrt{3}}\Bigr) + o(1),
  \qquad n \to \infty,
\]
where $\Phi$ is the standard normal cumulative distribution function
\[
  \Phi(z) \;=\; \frac{1}{\sqrt{2\pi}} \int_{-\infty}^{z}
                 e^{-t^{2}/2}\, \mathrm{d}t
\]
\citep[Thm.~1]{BayerDiaconis1992}, so that the approach to uniformity is
governed by the Gaussian tail $\int e^{-t^2/2}\,\mathrm{d}t$ evaluated at
the exponentially small argument $-2^{-\theta}/(4\sqrt 3)$. The scaling parameter
$\theta$ measures shuffles past the threshold $\tfrac{3}{2}\log_2 n$;
because the right-hand side collapses from $1$ to $0$ as $\theta$ runs
through an $O(1)$ window while the threshold itself grows like
$\log_2 n$, the transition is a true cutoff. For $n = 52$ this puts the
knee at $k \approx 8.55$ and yields distances dropping from $1.000$ at
$k=5$ to $0.334$ at $k=7$ and $0.167$ at $k=8$: the origin of ``seven
shuffles.'' The technical engine behind such exact statements is the
representation theory of finite groups: for the walk generated by random
transpositions, \citet{DiaconisShahshahani1981} proved cutoff at
$\tfrac12 n \log n$ by Fourier analysis on $S_n$, with the sharp
exponential rate
$\| Q^{\ast k} - U \|_{\mathrm{TV}} \le \tfrac12\, e^{-2c}$ once
$k = \tfrac12 n(\log n + c)$---the irreducible representations, indexed
by Young diagrams, diagonalize the walk, and a single upper-bound lemma
converts character sums into sharp convergence inequalities. The method is expounded in Diaconis's
monograph~\citep{Diaconis1988}, and it is the standing model of what a
\emph{hard convergence theorem} for an iterated mixing process looks
like: not ``eventually close to uniform,'' but \emph{exactly when},
with matching bounds.

\paragraph{Random walks on the hypercube, and the cutoff phenomenon.}
Between the card-shuffling summit and our object sits the solved mixing
problem closest to \SHA{}'s own geometry: the random walk on the
hypercube $\Ftwo^n$---the very space of \cref{sec:statespace}, walked
one coordinate at a time. At each step, pick one of the $n$ bit
positions uniformly and replace that bit by a fresh coin flip (this is
the Ehrenfest urn of statistical mechanics, wearing bits). How many
steps until the position is uniform on all $2^n$ points? The answer is
a second exact cutoff theorem: the total-variation distance stays
pinned near $1$ until $\tfrac14\, n \ln n$ steps, then collapses to $0$
inside a window of width $O(n)$---asymptotically instantaneous
(\cref{fig:cutoff}; the phenomenon, its history, and this example are
surveyed in~\citealp{Diaconis1996}). The threshold is pure coupon
collecting: a coordinate never yet touched is a coordinate whose bit is
still the initial one, and $\tfrac14 n \ln n$ is when the last
stragglers have been visited; the logarithm is the price of the
\emph{slowest} coordinate, not the average one---a moral worth
remembering whenever ``rounds needed to mix'' is discussed.

Better still, the machinery is transparent enough to expose a
coincidence this paper finds genuinely striking. On the abelian group
$\Ftwo^n$ the Fourier analysis that diagonalized card shuffles becomes
elementary: the characters are the Walsh functions
$x \mapsto (-1)^{\langle u, x\rangle}$, every translation-invariant
walk has them as eigenvectors, and \emph{the eigenvalues are exactly
the Walsh--Hadamard coefficients of the one-step distribution}. Mixing
is slow precisely in the directions $u$ where that coefficient is
large. Now reread the empirical ledger of \cref{sec:below}: ``flat
Walsh--Hadamard spectrum, no good linear approximation'' is the
property \emph{linear cryptanalysis} probes. The quantity the
cryptanalyst hunts and the quantity that governs mixing are the same
number: a linear distinguisher is a slow eigenvector, caught in the
act. Mixing analysis and linear cryptanalysis are, in this abelian
setting, one computation performed by two communities that rarely cite
each other---and the projects of \cref{sec:gap} amount to measuring
\SHA{}'s spectral gap experimentally, round by round. For scale: at
$n = 256$ the theorem's threshold is $\tfrac14\, n \ln n \approx 355$
single-coordinate updates, while one \SHA{} round nominally touches all
$256$ state bits at once---but in correlated fashion, which is exactly
why no theorem applies and the measured curve is the interesting
object.

\begin{figure}[ht]
  \centering
  \begin{tikzpicture}[font=\small, xscale=3.1, yscale=2.3]
    % axes
    \draw[->] (0,0) -- (2.25,0) node[below left=0.5mm and -6mm]
      {steps $\div$ $\tfrac14 n \ln n$};
    \draw[->] (0,0) -- (0,1.18) node[rotate=90, above left=1mm and 3mm]
      {distance to uniform};
    \draw (0.03,1) -- (-0.03,1) node[left, font=\scriptsize] {$1$};
    \draw (1,0.03) -- (1,-0.03) node[below, font=\scriptsize] {$1$};
    \draw[dashed, gray] (1,0) -- (1,1.05);
    % cutoff curves, sharpening with n (precomputed: 0.5*(1-tanh(a(x-1))))
    \draw[thick, msgfill!60!black] plot[smooth] coordinates {
      (0.00,0.982) (0.20,0.961) (0.40,0.917) (0.60,0.832) (0.80,0.690)
      (1.00,0.500) (1.20,0.310) (1.40,0.168) (1.60,0.083) (1.80,0.039)
      (2.00,0.018) (2.20,0.008)};
    \draw[thick, projframe] plot[smooth] coordinates {
      (0.00,1.000) (0.20,0.998) (0.40,0.992) (0.60,0.961) (0.70,0.917)
      (0.80,0.832) (0.90,0.690) (1.00,0.500) (1.10,0.310) (1.20,0.168)
      (1.30,0.083) (1.50,0.018) (1.70,0.004) (2.00,0.000) (2.20,0.000)};
    \draw[very thick, citewine] plot[smooth] coordinates {
      (0.00,1.000) (0.40,1.000) (0.60,0.999) (0.70,0.996) (0.80,0.973)
      (0.90,0.858) (1.00,0.500) (1.10,0.142) (1.20,0.027) (1.30,0.004)
      (1.50,0.000) (1.80,0.000) (2.20,0.000)};
    \node[font=\scriptsize, msgfill!60!black, anchor=west] at (1.32,0.30)
      {$n$ small};
    \node[font=\scriptsize, citewine, anchor=west] at (1.13,0.10)
      {$n$ large};
    \node[font=\scriptsize, gray, anchor=south, align=center]
      at (1.72,0.62) {window shrinks\\ relative to threshold};
    \draw[->, gray, font=\scriptsize] (1.45,0.60) -- (1.12,0.45);
  \end{tikzpicture}
  \caption{The cutoff phenomenon, schematically: distance to
  uniformity for the hypercube walk stays near $1$, then collapses
  within a window that is asymptotically negligible compared to the
  $\tfrac14 n \ln n$ threshold. ``How many rounds
  are enough?'' has, for this toy relative of \SHA{}, an answer sharp
  to within lower-order terms---and the toy-compression-walk project
  of this section asks whether \SHA{}-family walks show the same
  knee.}
  \label{fig:cutoff}
\end{figure}

The bridge to our object is not an analogy but a change of alphabet.
Every keyed block cipher is a family of permutations, and \SHA{}'s own
compression function is one in disguise (its core is a block cipher,
run in Davies--Meyer mode, \cref{sec:roundshape}). More directly: the
chaining recursion $h_i = f(h_{i-1}, m_i)$ of \cref{sec:merkledamgard},
fed \emph{random} message blocks, is literally a random walk on the
state space $\State$, driven by $2^{512}$ generators. Diaconis's
questions transfer verbatim---how many blocks until the distribution of
$h_i$ is uniform? is there cutoff? what is the spectral gap?---and, to
our knowledge, none has been answered, or seriously asked, for any
member of the SHA family at any scale. The obstruction is instructive:
the representation-theoretic engine wants a group acting on itself, and
the walk lives merely on a set---unless one exploits the two group
structures the set actually carries (\cref{sec:statespace}). Here is a
place where the paper's finite geometry and a fully developed classical
theory face each other across a gap that small experiments can begin to
survey.

\begin{project}{Cutoff for toy compression walks}{probability; linear
  algebra; programming}{4--6 weeks}
  Shrink the compression function (\cref{sec:roundshape}) to an $8$- or
  $16$-bit state. Random message blocks then make chaining a random
  walk on $2^{8}$ or $2^{16}$ points: build its transition matrix,
  compute the spectral gap, plot total-variation distance to uniformity
  against the number of blocks, and look for a cutoff profile. Vary the
  rounds-per-block and watch the gap and the cutoff window move. Every
  quantity here is exactly computable at toy scale, and the questions
  are verbatim those of the shuffling literature.
\end{project}

The interlude's moral for \cref{sec:gap}: on one flank of \SHA{} stands
a complete algebraic theory of linear generators, on the other a
complete probabilistic theory of mixing walks on groups; the function
itself was engineered into the mathematical no-man's-land between
them---touching complex analysis through \cref{sec:randommaps}, finite
fields and $2$-adic analysis through \cref{sec:statespace},
representation theory through the walk just described, and complexity
theory through \cref{sec:definitions}---and nearly every one of those
contacts is an unfollowed lead rather than a developed theory.

\subsection{The gap, stated plainly}\label{sec:gap}

Now place the object among its relatives, the pseudorandom number
generators, along two informal axes: how fast it is, and how much theory
supports its randomness claims (\cref{fig:prngspace}).

\begin{figure}[ht]
  \centering
  \begin{tikzpicture}[scale=1.0, font=\small]
    \draw[->] (0,0) -- (8.8,0) node[below left=1mm and 0mm]
      {analytical understanding};
    \draw[->] (0,0) -- (0,4.8) node[rotate=90, above left=1mm and 14mm]
      {speed / practicality};
    % fully understood, weak
    \node[align=center] (lcg) at (7.1,4.2)
      {LCG\\ \scriptsize lattice structure: solved, weak};
    \node[align=center] (mt) at (6.3,3.1)
      {Mersenne Twister\\ \scriptsize $\Ftwo$-linear: solved, distinguishable};
    % provable, slow
    \node[align=center] (bbs) at (6.7,0.9)
      {number-theoretic PRGs\\ \scriptsize provable from factoring-type\\
       \scriptsize assumptions, slow};
    % SHA
    \node[draw, very thick, circle, inner sep=1.5pt, fill=projfill,
          align=center] (sha) at (1.2,4.05) {\SHA};
    \node[gray, font=\scriptsize, align=left, anchor=west]
      at (1.95,4.05) {passes every test,\\ no theory at all};
    % the gap: brace bulging toward the empty lower-left region
    \draw[decorate, decoration={brace, mirror, amplitude=6pt}]
      (2.2,3.5) -- (5.0,1.55);
    \node[rotate=-35, font=\footnotesize] at (2.75,1.7)
      {\emph{the gap}: no intermediate landmarks};
  \end{tikzpicture}
  \caption{A cartoon of PRNG space. Classical generators are fast and
  fully understood---which is exactly how one knows they are weak.
  Number-theoretic generators carry proofs (conditional on hardness
  assumptions) and are slow. \SHA{} sits alone: fastest, empirically
  strongest, and analytically blank.}
  \label{fig:prngspace}
\end{figure}

The classical generators are \emph{transparent}: linear congruential
generators are so well understood that their defects (Marsaglia's
lattice planes) are theorems; twisters are $\Ftwo$-linear and hence
completely analyzable, and completely distinguishable. At the other
extreme, number-theoretic generators (of Blum--Micali type) come with
genuine reductions---distinguishing their output is provably as hard as
factoring-type problems---purchased at a price of thousands of arithmetic
operations per output bit. \SHA{} belongs to neither culture: it has
\emph{no} reduction to any studied hard problem, \emph{no} algebraic
model that survives even a handful of rounds, and \emph{no} distinguisher
after twenty-five years. It is simultaneously the best-tested and least
theorized object in the entire space.

Why should the gap be so large? Part of the answer is a legitimate
excuse: an unconditional proof that \SHA{} (or anything) is a one-way or
pseudorandom function would imply $\mathrm{P} \ne \mathrm{NP}$ and more;
in Impagliazzo's celebrated taxonomy of possible computational worlds,
whether we live in a world with one-way functions at all
(``Minicrypt''/``Cryptomania'' versus ``Algorithmica''/``Pessiland'')
is precisely the great open question~\citep{Impagliazzo1995}. Nobody
should demand a proof. But the excuse covers less than the gap. Between
``unconditional proof'' and ``nothing'' there is an enormous unexplored
middle: reductions to \emph{any} clean assumption; partial structure
theorems for few rounds; quantitative mixing results in either of the two
geometries of \cref{sec:statespace}; distributional results in Flajolet's
coordinates for scaled-down families (\cref{sec:randommaps}). For linear
generators that middle territory is dense with theorems; for ARX
functions it is nearly empty---not, as far as we can tell, because
theorems are known to be impossible, but because the object fell into a
disciplinary crack: too applied for one community, too structureless for
another, and (a hypothesis this prospectus exists to test) never yet
attacked with the analytic-combinatorics instrument at the appropriate
model-organism scale.

That is the research program: not ``prove \SHA{} secure''---nobody sane
proposes that---but \emph{populate the gap with measurements}: an atlas
of scaled-down \SHA{}-family maps (fewer rounds, narrower words, fewer
words), located in Flajolet's coordinate system, with the two-geometry
anatomy of \cref{sec:statespace} guiding which scaled families to
compare. Every panel of that atlas is a small, finishable, publishable
piece of natural history; \cref{sec:projects} cuts it into
undergraduate-sized pieces.

\begin{project}{Where do the tests start failing?}{programming;
  statistics}{4--6 weeks}
  Reduced-round \SHA{} must fail statistical batteries---at one round it
  is nearly raw message. Find the phase transition: run
  SP\,800-22~\citep{NIST2010} and TestU01~\citep{LEcuyerSimard2007}
  against generators built from $r$-round \SHA{} for $r = 1, 2, \dots$
  and chart which tests fail at which $r$. Compare with the round counts
  where cryptanalysis stalls (\cref{sec:below}). The two ``difficulty
  frontiers'' have, to our knowledge, never been plotted on one axis.
\end{project}

\begin{project}{Avalanche cartography}{programming; linear
  algebra}{4--6 weeks}
  Flip one input bit; which of the 256 output bits flip? For the true
  random-function silhouette each flips independently with probability
  $\tfrac12$~\citep{WebsterTavares1985}. Measure the full $512 \times 256$
  influence matrix round by round, render it as heatmaps, and quantify
  convergence to $\tfrac12$ (in what norm? after how many rounds?
  uniformly over bit positions?). Then close the loop to the ledger of
  \cref{sec:below}: confirm that full-round avalanche forces the counter
  stream's lag-$1$ autocorrelation to zero, and measure how many rounds
  are needed before that self-correlation vanishes. Relate the observed
  diffusion speed to the rotation constants of
  \cref{sec:roundshape}---the designers' choices become visible as
  geography.
\end{project}

\subsection{The road not taken: hashes with theorems}\label{sec:provable}

One region of \cref{fig:prngspace} deserves its own visit before we
leave the map: the lower right, where proofs live and speed dies. It is
not empty of hash functions. \emph{Provably secure} hash functions
exist, several are beautiful mathematics, and the world declined every
one of them; both halves of that sentence are instructive.

Exhibit one: VSH, ``very smooth hash''~\citep{ContiniLenstraSteinfeld2006}.
Its collision resistance is an honest theorem, conditional on a single
clean number-theoretic assumption (finding congruences of a certain
smooth-square kind modulo a hard-to-factor $n$), and it is fast by the
standards of its genre---only $25$ times or so slower than the SHA
family, where earlier provable designs had been slower by thousands.
But the same multiplicative structure that powers the proof is
\emph{visible in the outputs}: VSH is not remotely a random-oracle
stand-in (short relations among hash values exist by construction), so
of the whole catalogue of \cref{sec:uses} it can serve collision-bound
uses only. The proof purchases one property by installing exactly the
kind of algebraic transparency that \cref{sec:prngs} showed to be fatal
to all the others.

Exhibit two is closer to this paper's geometry. \emph{Cayley hashes}
walk a graph: fix a group $G$ and two generators $g_0, g_1$; hash a bit
string by reading it as directions,
$m_1 m_2 \cdots m_k \mapsto g_{m_1} g_{m_2} \cdots g_{m_k}$. Tillich
and Z{\'e}mor proposed $G = \mathrm{SL}_2(\mathbb{F}_{2^n})$ in
1994~\citep{TillichZemor1994}; Charles, Goren, and Lauter modernized
the idea with Lubotzky--Phillips--Sarnak Ramanujan graphs and, in a
second variant, supersingular isogeny graphs of elliptic
curves~\citep{CharlesGorenLauter2009}. Two theorems come for free.
A collision is a pair of distinct words with the same product---so
collision resistance \emph{is} a group-theoretic statement (no short
relations among the generators are findable), inherited whole from a
recognizable mathematical problem. And output quality is spectral: if
the Cayley graph is an \emph{expander}---nontrivial eigenvalues
uniformly bounded away from the degree---then the distribution of
hashes of length-$k$ messages converges to uniform at an exponential
rate governed by the spectral gap~\citep{HooryLinialWigderson2006}.
The mixing mathematics of \cref{sec:prngs} returns, no longer as
analysis but as a \emph{design principle}: these are hash functions
whose avalanche property is a theorem about eigenvalues.

Then the pothole. In 2008 the LPS instantiation was broken---collisions
computed outright---by lifting the problem from the finite group to a
characteristic-zero arithmetic where Euclidean-style algorithms find
short relations; the authors of the break were Tillich and Z{\'e}mor,
the designers of 1994, and their own original $\mathrm{SL}_2$ proposal
later fell, for its published parameters, to a similar
lift~\citep{TillichZemor2008}. Note carefully what died. The theorems
were and remain true; the \emph{assumptions} were false---the
underlying group problems, young and algebraically rich, simply were
not hard. (The isogeny-graph sibling's assumption has so far survived,
and seeded a whole branch of post-quantum cryptography.) The moral is
the sober one: a security proof does not abolish risk, it relocates
risk from the object to the assumption, and an assumption chosen for
its provability can be weaker than ignorance. Set beside
\cref{sec:falling}: MD5 fell \emph{with} twenty years of massive use
and no theorem; LPS fell \emph{with} a theorem and almost no use. The
fossil record refuses to pick a side.

Why, then, does the world run on the theorem-free object? Speed is the
visible reason---two orders of magnitude---but the load-bearing one is
now easy to state with this section's vocabulary: every use of
\cref{sec:uses} beyond bare collision resistance needs the
random-oracle \emph{pretension}, behaving like structurelessness
itself, and a provable design is structured \emph{by construction};
its proof and its disqualification are the same fact. The market
equilibrium of \cref{fig:prngspace} is thus not an accident of
laziness: given the choice between an object about which something is
proved and something else is visibly false, and an object about which
nothing is proved and nothing is visibly false, every deployed system
chose the second, with Rogaway's human-ignorance calculus
(\cref{sec:definitions}) as the honest license. The gap is not a hole
in the market; it is the market.

\begin{project}{Cayley hashes by hand}{group theory; linear algebra;
  programming}{3--4 weeks}
  Implement the Tillich--Z{\'e}mor hash~\citep{TillichZemor1994} at toy
  scale: $\mathrm{SL}_2(\mathbb{F}_{2^n})$ for $n = 8, \dots, 16$, two
  generator matrices, hash $=$ matrix product read off the message
  bits. (i) Measure output equidistribution against message length and
  match the exponential rate to the spectral gap of the Cayley graph,
  computable by eigenvalue methods at this
  scale~\citep{HooryLinialWigderson2006}---the mixing theory of
  \cref{sec:prngs} as a bench experiment. (ii) Hunt collisions two
  ways, generic birthday search versus structure (palindromic words,
  elements of small order, subfield membership), and compare costs.
  (iii) As a stretch goal, read the lifting attack~\citep{TillichZemor2008}
  and implement its Euclidean core for tiny parameters. Deliverable: a
  working Cayley hash, gap-versus-equidistribution plots, and a
  collision gallery---provable cryptography's rise and fall on one
  laptop.
\end{project}
